%=============================
%  SEGUNDO CAPITULO
%=============================
\chapter{FUNDAMENTOS Y MÉTRICAS DE AGUJEROS NEGROS EN EL SISTEMA INTERNACIONAL}

\section{Ecuaciones de Einstein-Maxwell }

En el Sistema Internacional de Unidades (SI), la dinámica acoplada entre la geometría espaciotemporal tetradimensional $\mathcal{M}$ y el campo electromagnético está gobernada por la Acción de Einstein-Maxwell \citep{M_P_Hobson_General_Relativity, Sean_Carroll_Spacetime, 05villegas}:
\begin{equation}
    S[g^{\mu\nu}, A_\alpha] = \int_{\mathcal{M}} \left( \frac{c^4}{16\pi G} R - \frac{1}{4\mu_0} F_{\alpha\beta}F^{\alpha\beta} \right) \sqrt{-g} \, d^4x + S_m
    \label{eq:accion_EM_SI}
\end{equation}
donde $G$ es la constante de gravitación universal, $c$ es la velocidad de la luz en el vacío, $\mu_0$ es la permeabilidad magnética del vacío, $g \equiv \det(g_{\mu\nu}) < 0$ es el determinante de la métrica, $R = g^{\mu\nu} R_{\mu\nu}$ es el escalar de Ricci, y $F_{\alpha\beta} \equiv \partial_\alpha A_\beta - \partial_\beta A_\alpha$ es el tensor electromagnético expresado en términos del cuadripotencial $A_\alpha$. El elemento de línea plano en coordenadas esféricas adopta la forma $ds^2 = c^2 dt^2 - dr^2 - r^2 d\theta^2 - r^2 \sin^2\theta d\phi^2$ \citep{M_P_Hobson_General_Relativity, 05villegas}.

Para obtener las ecuaciones de campo de Einstein-Maxwell, aplicamos el principio de mínima acción variando $S$ respecto a la métrica contravariante $g^{\mu\nu}$, manteniendo fijo el campo electromagnético $A_\alpha$. En primer lugar, utilizando la identidad de Jacobi $\delta g = g \, g^{\mu\nu} \delta g_{\mu\nu} = -g \, g_{\mu\nu} \delta g^{\mu\nu}$, la variación del elemento de volumen $\sqrt{-g}$ resulta:
\begin{equation}
    \delta \sqrt{-g} = -\frac{1}{2\sqrt{-g}} \delta g
    \label{eq:var_vol_1}
\end{equation}
\begin{equation}
    \delta \sqrt{-g} = -\frac{1}{2} \sqrt{-g} \, g_{\mu\nu} \delta g^{\mu\nu}
    \label{eq:var_vol_2}
\end{equation}

A continuación, calculamos la variación de la parte gravitacional de la acción, $S_{\mathrm{EH}} = \frac{c^4}{16\pi G} \int R \sqrt{-g} \, d^4x$:
\begin{equation}
    \delta S_{\mathrm{EH}} = \frac{c^4}{16\pi G} \int \left( R_{\mu\nu} \delta g^{\mu\nu} \sqrt{-g} + g^{\mu\nu} \delta R_{\mu\nu} \sqrt{-g} + R \, \delta \sqrt{-g} \right) d^4x
    \label{eq:var_seh_1}
\end{equation}
Por la identidad de Palatini, el término $g^{\mu\nu} \delta R_{\mu\nu} = \nabla_\alpha (\nabla_\beta \delta g^{\alpha\beta} - g^{\alpha\beta} \nabla_\gamma \delta g^{\gamma}_{\phantom{\gamma}\alpha})$ constituye una divergencia total que se anula en la frontera del espacio-tiempo mediante el teorema de Gauss. Sustituyendo la ec.~(\ref{eq:var_vol_2}) en la ec.~(\ref{eq:var_seh_1}), la variación del término de Einstein-Hilbert toma la forma:
\begin{equation}
    \delta S_{\mathrm{EH}} = \frac{c^4}{16\pi G} \int \left( R_{\mu\nu} - \frac{1}{2} g_{\mu\nu} R \right) \delta g^{\mu\nu} \sqrt{-g} \, d^4x
    \label{eq:var_seh_2}
\end{equation}

Por otro lado, para variar el término electromagnético expresamos la cantidad escalar como $F_{\alpha\beta} F^{\alpha\beta} = g^{\alpha\gamma} g^{\beta\delta} F_{\alpha\beta} F_{\gamma\delta}$. La variación directa respecto a la métrica $g^{\mu\nu}$ da:
\begin{equation}
    \delta \left( F_{\alpha\beta} F^{\alpha\beta} \right) = \delta(g^{\alpha\gamma}) g^{\beta\delta} F_{\alpha\beta} F_{\gamma\delta} + g^{\alpha\gamma} \delta(g^{\beta\delta}) F_{\alpha\beta} F_{\gamma\delta}
    \label{eq:var_f2_1}
\end{equation}
\begin{equation}
    \delta \left( F_{\alpha\beta} F^{\alpha\beta} \right) = 2 F_{\mu\alpha} F_{\nu}^{\phantom{\nu}\alpha} \, \delta g^{\mu\nu}
    \label{eq:var_f2_2}
\end{equation}
Con este resultado, la variación de la acción electromagnética $S_{\mathrm{EM}} = -\frac{1}{4\mu_0} \int F_{\alpha\beta} F^{\alpha\beta} \sqrt{-g} \, d^4x$ se escribe como:
\begin{equation}
    \delta S_{\mathrm{EM}} = -\frac{1}{4\mu_0} \int \left[ \delta\left(F_{\alpha\beta} F^{\alpha\beta}\right) \sqrt{-g} + F_{\alpha\beta} F^{\alpha\beta} \delta\sqrt{-g} \right] d^4x
    \label{eq:var_sem_1}
\end{equation}
Sustituyendo las ecs.~(\ref{eq:var_vol_2}) y~(\ref{eq:var_f2_2}) en la ec.~(\ref{eq:var_sem_1}):
\begin{equation}
    \delta S_{\mathrm{EM}} = -\frac{1}{4\mu_0} \int \left[ 2 F_{\mu\alpha} F_{\nu}^{\phantom{\nu}\alpha} - \frac{1}{2} g_{\mu\nu} F_{\alpha\beta} F^{\alpha\beta} \right] \delta g^{\mu\nu} \sqrt{-g} \, d^4x
    \label{eq:var_sem_2}
\end{equation}
\begin{equation}
    \delta S_{\mathrm{EM}} = -\frac{1}{2} \int \left[ \frac{1}{\mu_0} \left( F_{\mu\,\alpha}F_\nu^{\phantom{\nu}\alpha} - \frac{1}{4}g_{\mu\nu} F_{\alpha\beta}F^{\alpha\beta} \right) \right] \delta g^{\mu\nu} \sqrt{-g} \, d^4x
    \label{eq:var_sem_3}
\end{equation}
Esta última expresión nos permite definir el tensor de energía-impulso electromagnético en el SI como \citep{M_P_Hobson_General_Relativity, 05villegas}:
\begin{equation}
    T_{\mu\nu}^{(\mathrm{EM})} \equiv -\frac{2}{\sqrt{-g}} \frac{\delta S_{\mathrm{EM}}}{\delta g^{\mu\nu}}
    \label{eq:def_tensor_EM}
\end{equation}
\begin{equation}
    T_{\mu\nu}^{(\mathrm{EM})} = \frac{1}{\mu_0} \left( F_{\mu\,\alpha}F_\nu^{\phantom{\nu}\alpha} - \frac{1}{4}g_{\mu\nu} F_{\alpha\beta}F^{\alpha\beta} \right)
    \label{eq:tensor_EM_SI}
\end{equation}

Exigiendo la condición de punto estacionario $\delta S = \delta S_{\mathrm{EH}} + \delta S_{\mathrm{EM}} = 0$ para cualquier variación arbitraria $\delta g^{\mu\nu}$, obtenemos las ecuaciones de campo de Einstein en el SI \citep{M_P_Hobson_General_Relativity, Sean_Carroll_Spacetime}:
\begin{equation}
    R_{\mu\nu} - \frac{1}{2}g_{\mu\nu}R = \frac{8\pi G}{c^4} T_{\mu\nu}^{(\mathrm{EM})}
    \label{eq:einstein_SI}
\end{equation}
Para simplificar esta ecuación, calculamos la traza del tensor de energía-impulso $T^{(\mathrm{EM})} \equiv g^{\mu\nu} T_{\mu\nu}^{(\mathrm{EM})}$ en cuatro dimensiones:
\begin{equation}
    T^{(\mathrm{EM})} = \frac{1}{\mu_0} \left( g^{\mu\nu} F_{\mu\alpha} F_{\nu}^{\phantom{\nu}\alpha} - \frac{1}{4} g^{\mu\nu} g_{\mu\nu} F_{\alpha\beta} F^{\alpha\beta} \right)
    \label{eq:traza_tem_1}
\end{equation}
Dado que $g^{\mu\nu} g_{\mu\nu} = 4$, la ec.~(\ref{eq:traza_tem_1}) se reduce a:
\begin{equation}
    T^{(\mathrm{EM})} = \frac{1}{\mu_0} \left( F_{\alpha\beta} F^{\alpha\beta} - \frac{4}{4} F_{\alpha\beta} F^{\alpha\beta} \right) \equiv 0
    \label{eq:traza_tem_2}
\end{equation}
Al ser $T^{(\mathrm{EM})} = 0$, contratando la ec.~(\ref{eq:einstein_SI}) con $g^{\mu\nu}$ obtenemos $R - \frac{1}{2}(4)R = 0 \implies -R = 0 \implies R = 0$. Por lo tanto, en el electrovacío el escalar de Ricci se anula idénticamente y las ecuaciones de Einstein adoptan la forma simplificada:
\begin{equation}
    R_{\mu\nu} = \frac{8\pi G}{c^4} T_{\mu\nu}^{(\mathrm{EM})}
    \label{eq:einstein_electrovacio}
\end{equation}

Finalmente, derivamos las ecuaciones de Maxwell variando la acción $S_{\mathrm{EM}}$ respecto al cuadripotencial $A_\beta$. Utilizando la relación $\delta F_{\alpha\beta} = \partial_\alpha \delta A_\beta - \partial_\beta \delta A_\alpha = 2 \nabla_{[\alpha} \delta A_{\beta]}$:
\begin{equation}
    \delta_A S_{\mathrm{EM}} = -\frac{1}{\mu_0} \int F^{\alpha\beta} (\nabla_\alpha \delta A_\beta) \sqrt{-g} \, d^4x
    \label{eq:var_maxwell_1}
\end{equation}
Integrando por partes y aplicando la compatibilidad métrica $\nabla_\alpha \sqrt{-g} = 0$:
\begin{equation}
    \delta_A S_{\mathrm{EM}} = \frac{1}{\mu_0} \int \left( \nabla_\alpha F^{\alpha\beta} \right) \delta A_\beta \sqrt{-g} \, d^4x = 0
    \label{eq:var_maxwell_2}
\end{equation}
De la arbitrariedad de $\delta A_\beta$ se derivan las ecuaciones de Maxwell covariantes en el vacío \citep{M_P_Hobson_General_Relativity, Sean_Carroll_Spacetime}:
\begin{equation}
    \nabla_\alpha F^{\alpha\beta} = 0
    \label{eq:maxwell_covariante}
\end{equation}
que en términos de derivadas parciales se expresan como:
\begin{equation}
    \frac{1}{\sqrt{-g}} \partial_\alpha \left( \sqrt{-g} \, F^{\alpha\beta} \right) = 0
    \label{eq:maxwell_parciales}
\end{equation}
junto con la identidad de Bianchi electromagnética:
\begin{equation}
    \partial_\alpha F_{\beta\gamma} + \partial_\beta F_{\gamma\alpha} + \partial_\gamma F_{\alpha\beta} = 0
    \label{eq:bianchi_em}
\end{equation}
El conjunto acoplado de las ecuaciones~(\ref{eq:einstein_electrovacio}), (\ref{eq:maxwell_parciales}) y~(\ref{eq:bianchi_em}) constituye el sistema de Einstein-Maxwell  cuya solución exacta e incompresible para un agujero negro con masa $M$, momento angular $J$ y carga $Q$ es la métrica de Kerr-Newman \citep{1403.5234v2, 05villegas}.




\section{La métrica de Kerr-Newman en unidades SI}

Al incorporar el momento angular $J$ del objeto compacto, la rotación destruye la simetría esférica pero conserva la simetría axial respecto al eje de rotación y la estacionariedad temporal \citep{1403.5234v2, Sean_Carroll_Spacetime}. El parámetro de espín con dimensiones de longitud se define en el SI como \citep{M_P_Hobson_General_Relativity, s10052-021-09180-y}:
\begin{equation}
    a \equiv \frac{J}{M c}
    \label{eq:parametro_espin_SI}
\end{equation}

En coordenadas de Boyer-Lindquist $(t, r, \theta, \phi)$ y bajo la convención de firma de Minkowski $(+,-,-,-)$, el elemento de línea exacto del espaciotiempo de Kerr-Newman se expresa en el SI como \citep{1403.5234v2, pdfcoffee.com_chandrasekhar}:
\begin{equation}
    ds^2 = g_{tt} c^2 dt^2 + 2 c \, g_{t\phi} dt d\phi + g_{rr} dr^2 + g_{\theta\theta} d\theta^2 + g_{\phi\phi} d\phi^2
    \label{eq:linea_kerr_newman_general}
\end{equation}
donde las componentes covariantes del tensor métrico $g_{\mu\nu}$ son \citep{1403.5234v2, s10052-021-09180-y}:
\begin{equation}
    g_{tt}(r,\theta) = 1 - \frac{2r_g r - r_g^2 \delta}{\Sigma(r,\theta)}
    \label{eq:gtt_KN}
\end{equation}
\begin{equation}
    g_{t\phi}(r,\theta) = \frac{a \sin^2\theta}{\Sigma(r,\theta)} \left( 2r_g r - r_g^2 \delta \right)
    \label{eq:gtphi_KN}
\end{equation}
\begin{equation}
    g_{rr}(r,\theta) = -\frac{\Sigma(r,\theta)}{\Delta(r)}
    \label{eq:grr_KN}
\end{equation}
\begin{equation}
    g_{\theta\theta}(r,\theta) = -\Sigma(r,\theta)
    \label{eq:gthetatheta_KN}
\end{equation}
\begin{equation}
    g_{\phi\phi}(r,\theta) = -\left[ r^2 + a^2 + \frac{a^2 \sin^2\theta}{\Sigma(r,\theta)} \left( 2r_g r - r_g^2 \delta \right) \right] \sin^2\theta
    \label{eq:gphiphi_KN}
\end{equation}
Las funciones métricas auxiliares $\Sigma(r,\theta)$ y $\Delta(r)$ se definen explícitamente como \citep{1403.5234v2, s10052-021-09180-y}:
\begin{equation}
    \Sigma(r,\theta) = r^2 + a^2 \cos^2\theta
    \label{eq:sigma_def_detallada}
\end{equation}
\begin{equation}
    \Delta(r) = r^2 - \frac{2GM}{c^2}r + a^2 + \frac{G Q^2}{4\pi\varepsilon_0 c^4} = r^2 - 2 r_g r + a^2 + r_g^2 \delta
    \label{eq:delta_def_detallada}
\end{equation}

El cuadripotencial electromagnético acoplado a la geometría de Kerr-Newman viene dado en coordenadas de Boyer-Lindquist por \citep{1403.5234v2, Sean_Carroll_Spacetime}:
\begin{equation}
    A_\mu = \frac{Q r}{4\pi\varepsilon_0 \Sigma c} \left( 1, 0, 0, -a \sin^2\theta \right)
    \label{eq:cuadripotencial_KN}
\end{equation}
donde la componente temporal $A_t = \frac{Q r}{4\pi\varepsilon_0 \Sigma c}$ representa el potencial eléctrico modificado por la rotación y $A_\phi = -\frac{Q r a \sin^2\theta}{4\pi\varepsilon_0 \Sigma c}$ describe el campo magnético dipolar inducido por el arrastre del momento angular \citep{M_P_Hobson_General_Relativity}.

El determinante de la métrica de Kerr-Newman se calcula directamente a partir del bloque no diagonal $(t,\phi)$ y los componentes diagonales $(r,\theta)$:
\begin{equation}
    g \equiv \det(g_{\mu\nu}) = (g_{tt} g_{\phi\phi} - g_{t\phi}^2) g_{rr} g_{\theta\theta} = -\Sigma^2 \sin^2\theta
    \label{eq:det_kerr_newman}
\end{equation}
\begin{equation}
    \sqrt{-g} = \Sigma \sin^2\theta \dots \implies \sqrt{-g} = \Sigma \sin\theta
    \label{eq:raiz_det_kerr_newman}
\end{equation}

Invertimos la matriz métrica covariante $g_{\mu\nu}$ para obtener los componentes contravariantes $g^{\mu\nu}$, fundamentales para la formulación del potencial radial efectivo mediante la ecuación de Hamilton-Jacobi en el Capítulo III \citep{pdfcoffee.com_chandrasekhar, 1403.5234v2}:
\begin{equation}
    g^{tt}(r,\theta) = \frac{(r^2 + a^2)^2 - a^2 \Delta(r) \sin^2\theta}{\Sigma(r,\theta) \, \Delta(r)}
    \label{eq:g_contra_tt}
\end{equation}
\begin{equation}
    g^{t\phi}(r,\theta) = \frac{a \left[ (r^2 + a^2) - \Delta(r) \right]}{\Sigma(r,\theta) \, \Delta(r)} = \frac{a \left( 2 r_g r - r_g^2 \delta \right)}{\Sigma(r,\theta) \, \Delta(r)}
    \label{eq:g_contra_tphi}
\end{equation}
\begin{equation}
    g^{rr}(r,\theta) = -\frac{\Delta(r)}{\Sigma(r,\theta)}
    \label{eq:g_contra_rr}
\end{equation}
\begin{equation}
    g^{\theta\theta}(r,\theta) = -\frac{1}{\Sigma(r,\theta)}
    \label{eq:g_contra_thetatheta}
\end{equation}
\begin{equation}
    g^{\phi\phi}(r,\theta) = -\frac{\Delta(r) - a^2 \sin^2\theta}{\Sigma(r,\theta) \, \Delta(r) \sin^2\theta}
    \label{eq:g_contra_phiphi}
\end{equation}

La estructura global de la métrica de Kerr-Newman posee dos superficies características fundamentales:
\begin{enumerate}
    \item \textbf{Horizontes de eventos ($r_\pm$):} Corresponden a las hipersuperficies nulas dadas por las raíces de la función $\Delta(r) = 0$ \citep{1403.5234v2, s10052-021-09180-y}:
    \begin{equation}
        r_\pm = r_g \left( 1 \pm \sqrt{1 - \frac{a^2}{r_g^2} - \delta} \right)
        \label{eq:horizontes_KN}
    \end{equation}
    \item \textbf{Ergosfera o Límite Estático ($r_{\text{ergo}}(\theta)$):} Corresponde al lugar geométrico donde el componente $g_{tt}(r,\theta) = 0$, lo que impide la existencia de observadores estáticos con cuadrivelocidad paralela al vector de Killing temporal \citep{M_P_Hobson_General_Relativity, Sean_Carroll_Spacetime}:
    \begin{equation}
        r_{\text{ergo}}(\theta) = r_g \left( 1 + \sqrt{1 - \frac{a^2 \cos^2\theta}{r_g^2} - \delta} \right)
        \label{eq:ergosfera_KN}
    \end{equation}
\end{enumerate}
La condición para la existencia de un agujero negro no extremall exige que el discriminante sea estrictamente positivo: $a^2/r_g^2 + \delta < 1$. En el régimen de carga débil ($\delta \ll 1$), la métrica de Kerr-Newman~(\ref{eq:linea_kerr_newman_general}) constituye la perturbación exacta adecuada para analizar el impacto de la carga electrostática sobre la dinámica de fotones sin alterar la integrabilidad matemática del espaciotiempo \citep{1403.5234v2, s10052-021-09180-y}.